% Generated draft: requires mathematical and source review.
\documentclass[11pt]{article}
\usepackage[margin=25mm]{geometry}
\usepackage{fontspec}
\setmainfont{Latin Modern Roman}
\usepackage{amsmath,amssymb,mathtools,unicode-math}
\setmathfont{Latin Modern Math}
\usepackage{enumitem,needspace,xurl,hyperref}
\setlength{\parindent}{0pt}
\setlength{\parskip}{5pt}
\setlength{\emergencystretch}{3em}
\newcommand{\N}{\mathbb{N}}
\newcommand{\Z}{\mathbb{Z}}
\newcommand{\Q}{\mathbb{Q}}
\newcommand{\R}{\mathbb{R}}
\newcommand{\C}{\mathbb{C}}
\newcommand{\F}{\mathbb{F}}
\newcommand{\A}{\mathbb{A}}
\newcommand{\PP}{\mathbb{P}}
\newcommand{\E}{\mathbb{E}}
\newcommand{\eps}{\varepsilon}
\newcommand{\dd}{\,\mathrm{d}}
\begin{document}
\section*{UnsolvedMath ID: 30006457}
\textbf{Source problem code:} catalogue entry
\subsection{Definitions and mathematical statement}
Let \(q=p^r\), where \(p\) is an odd prime and \(r\geq1\) is an integer, and let \(\F_q\) be the field with \(q\) elements. For each integer \(g\geq0\), put
\[
\mathcal H_{2g+1}(q)=\{d\in\F_q[t]:d\text{ is monic and squarefree},\ \deg d=2g+1\}.
\]
Here squarefree means that no square of a nonconstant irreducible polynomial divides \(d\).

The supplied record does not define \(\chi_d\) or normalize its \(L\)-function. To make the statement precise, adopt the standard hyperelliptic convention as a provisional interpretation: for a monic irreducible polynomial \(P\), define \(\chi_d(P)=0\) if \(P\mid d\), and otherwise define \(\chi_d(P)\in\{1,-1\}\) according as the image of \(d\) is a square or a nonsquare in \((\F_q[t]/(P))^\times\). Extend multiplicatively to monic polynomials, with \(\chi_d(1)=1\). For \(d\in\mathcal H_{2g+1}(q)\), use the associated polynomial
\[
\mathcal L_d(u)=\sum_{n=0}^{2g}\left(\sum_{\substack{f\in\F_q[t]\text{ monic}\\\deg f=n}}\chi_d(f)\right)u^n,
\qquad L(s,\chi_d)=\mathcal L_d(q^{-s}).
\]
This convention must be checked against the source report; it is not explicitly specified in the supplied record.

The extracted conjecture is that, for every odd prime power \(q\) and every integer \(k\geq1\), there exists an explicit polynomial \(R_{k,q}(X)\), of degree exactly \(k(k+1)/2\), such that
\[
\lim_{\substack{g\to\infty\\g\in\Z_{\geq0}}}
\left(q^{-(2g+1)}\sum_{d\in\mathcal H_{2g+1}(q)}L(1/2,\chi_d)^k-R_{k,q}(g)\right)=0.
\]
Here \(q\) and \(k\) are fixed in the limit, and \(R_{k,q}\) is independent of \(g\). No uniformity in \(q\) or \(k\) is asserted. The normalization is exactly \(q^{-(2g+1)}\), rather than division by the cardinality of \(\mathcal H_{2g+1}(q)\). The supplied record gives neither a coefficient formula for \(R_{k,q}\) nor its coefficient field, and does not specify what form of explicit description is intended.

The source labels this statement Conjecture 3. Its accompanying discussion reports a proof when \(q\) is sufficiently large relative to \(k\), using homological stability and homotopy theory. It describes an error involving positive constants \(A,B\), but the extracted expression \(O(4g(k+1)q^{-Ag+B})\) has unresolved typography, particularly the possible exponent on \(4\). No precise error estimate or dependence of these constants is asserted here. The discussion says that sufficiently large \(q\) gives a power-saving error and permits recovery of every coefficient of the moment polynomials; it presents this as evidence for the Conrey--Farmer--Keating--Rubinstein--Snaith recipe. These are supplied source claims requiring review, not independently verified conclusions.
\subsection{Short English statement}
For each fixed odd prime power \(q\) and positive integer \(k\), the conjecture predicts that the \(k\)-th central moment of quadratic \(L\)-functions over monic squarefree polynomials of degree \(2g+1\), normalized by \(q^{2g+1}\), equals an explicit polynomial in \(g\) of degree exactly \(k(k+1)/2\), with an error tending to zero. The supplied curation labels the problem partially solved, reporting results for sufficiently large \(q\) relative to \(k\). The character convention and the extracted error estimate require source review.
\subsection{Sources}
\begin{itemize}
\item Supplied local UnsolvedMath registry record: permanent ID \(30006457\), source problem code \texttt{catalogue entry}, title \emph{Moments of Quadratic L-Functions over Function Fields}, in the Oberwolfach Reports Open Problems collection. No individual proposer is supplied; the proposed year is recorded as 2025. The registry's \texttt{corrected\_verified} statement label is accompanied by a note that only formula delimiters, notation, grammar, and typography were repaired; it does not constitute independent mathematical verification here.
\item Supplied report citation: \emph{Analytic Number Theory (2026)}, Oberwolfach Reports, pp.\ 2701--2764, DOI \url{https://doi.org/10.4171/owr/2025/51}. The excerpt labels the target Conjecture 3 and contains the page marker 2723. The supplied literature list calls this \emph{Oberwolfach Report 51/2025}. The differing year designations are retained without reconciliation.
\item Supplied literature reference: \emph{Hyperelliptic curves and moments}, \url{https://arxiv.org/abs/2302.07664}.
\item Supplied literature reference: \emph{Uniform twisted homological stability}, \url{https://arxiv.org/abs/2402.00354}.
\item Literature assessment dated exactly \texttt{2026-08-22}: the supplied curation labels the problem partially solved and reports that the function-field moment formula has been proved with an explicit error term for \(q\) sufficiently large relative to \(k\), while the formulation for every odd prime power and every \(k\) is not fully established. This is supplied dated curation derived from the OWR record, not an independently verified literature assessment. The supplied background expressly requests independent verification of time-sensitive status and cited claims.
\end{itemize}
\end{document}
